% !TEX program = xelatex
% Самодостаточный конспект. Компилировать XeLaTeX/LuaLaTeX (tectonic 11-stepennye-ryady.tex).
\documentclass[11pt,a4paper]{article}
\newcommand{\coursename}{Математический анализ}
\newcommand{\rightheader}{§11. Степенные ряды}
\newcommand{\doctitle}{Степенные ряды}
% ==== Общая преамбула конспектов (собирается в каждый .tex как самодостаточная) ====
% Перед этим файлом должны быть определены: \documentclass и макросы
%   \coursename  — название курса (левый колонтитул, подзаголовок),
%   \rightheader — правый колонтитул (напр. «§2. Рациональные дроби»),
%   \doctitle    — заголовок раздела на титуле.
\usepackage{fontspec}
\setmainfont{cmunrm.otf}[
  BoldFont       = cmunbx.otf,
  ItalicFont     = cmunti.otf,
  BoldItalicFont = cmunbi.otf]
\usepackage{polyglossia}
\setdefaultlanguage{russian}
\setotherlanguage{english}

\usepackage{amsmath,amssymb,amsthm,mathtools}
\usepackage[a4paper,margin=2.2cm]{geometry}
\usepackage{enumitem}
\usepackage{graphicx}
\usepackage{array}
\usepackage{xcolor}
\definecolor{accent}{HTML}{1F6FEB}
\definecolor{rulecol}{HTML}{D0D7DE}
\usepackage[colorlinks=true,linkcolor=accent,urlcolor=accent,citecolor=accent]{hyperref}
\usepackage{titlesec}
\usepackage{fancyhdr}

% ----- Колонтитулы -----
\pagestyle{fancy}
\fancyhf{}
\fancyhead[L]{\small\itshape \coursename}
\fancyhead[R]{\small\itshape \rightheader}
\fancyfoot[C]{\thepage}
\renewcommand{\headrulewidth}{0.4pt}
\renewcommand{\headrule}{\hbox to\headwidth{\color{rulecol}\leaders\hrule height \headrulewidth\hfill}}

% ----- Заголовки -----
\titleformat{\section}{\large\bfseries\color{accent}}{\thesection.}{0.6em}{}
\titleformat{\subsection}{\bfseries}{\thesubsection.}{0.5em}{}

% ----- Теоремные окружения (стиль конспекта) -----
\theoremstyle{definition}
\newtheorem{definition}{Определение}[section]
\newtheorem{example}[definition]{Пример}
\newtheorem{remark}[definition]{Замечание}
\theoremstyle{plain}
\newtheorem{theorem}[definition]{Теорема}
\newtheorem{lemma}[definition]{Лемма}
\newtheorem{corollary}[definition]{Следствие}
\newtheorem{property}[definition]{Свойство}
\newtheorem{criterion}[definition]{Критерий}
\renewcommand{\proofname}{Доказательство}

% ----- Операторы и сокращения -----
\DeclareMathOperator{\tg}{tg}
\DeclareMathOperator{\ctg}{ctg}
\DeclareMathOperator{\arctg}{arctg}
\DeclareMathOperator{\arcctg}{arcctg}
\DeclareMathOperator{\sh}{sh}
\DeclareMathOperator{\ch}{ch}
\DeclareMathOperator{\Th}{th}
\DeclareMathOperator{\cth}{cth}
\DeclareMathOperator{\Const}{const}
\DeclareMathOperator{\sign}{sign}
\DeclareMathOperator{\rang}{rang}
\DeclareMathOperator{\diam}{diam}
\DeclareMathOperator{\Span}{span}
\DeclareMathOperator{\Ker}{Ker}
\DeclareMathOperator{\Img}{Im}
\DeclareMathOperator{\tr}{tr}
\newcommand{\R}{\mathbb{R}}
\newcommand{\C}{\mathbb{C}}
\newcommand{\N}{\mathbb{N}}
\newcommand{\Z}{\mathbb{Z}}
\newcommand{\Q}{\mathbb{Q}}
\newcommand{\dx}{\,dx}
\newcommand{\dt}{\,dt}
\newcommand{\du}{\,du}
\newcommand{\eps}{\varepsilon}
\renewcommand{\Re}{\operatorname{Re}}
\renewcommand{\Im}{\operatorname{Im}}
\renewcommand{\le}{\leqslant}
\renewcommand{\ge}{\geqslant}
\newcommand{\scal}[2]{\left( #1,\,#2\right)}
\DeclarePairedDelimiter{\abs}{\lvert}{\rvert}
\DeclarePairedDelimiter{\norm}{\lVert}{\rVert}

\begin{document}
\begin{center}
  {\LARGE\bfseries \doctitle}\\[4pt]
  {\large\itshape \coursename}
\end{center}
\vspace{6pt}

\section{Степенной ряд и область сходимости}

\begin{definition}
\emph{Степенным рядом} называют функциональный ряд вида
\[
  \sum_{n=0}^\infty a_n(x-x_0)^n=a_0+a_1(x-x_0)+a_2(x-x_0)^2+\dots,
\]
где $a_n$ — коэффициенты, $x_0$ — центр ряда. Заменой $x-x_0\mapsto x$ всегда можно считать
$x_0=0$: $\sum a_n x^n$.
\end{definition}

\begin{theorem}[первая теорема Абеля]
Если ряд $\sum a_n x^n$ сходится в точке $x_0\ne0$, то он сходится абсолютно при всех $x$ с
$\abs x<\abs{x_0}$. Если ряд расходится в точке $x_1$, то он расходится при всех $\abs x>\abs{x_1}$.
\end{theorem}

\begin{proof}
Из сходимости $\sum a_n x_0^n$ следует $a_n x_0^n\to0$, значит $\abs{a_n x_0^n}\le M$. При
$\abs x<\abs{x_0}$, положив $q=\abs{x/x_0}<1$,
\[
  \abs{a_n x^n}=\abs{a_n x_0^n}\cdot\abs[\Big]{\frac{x}{x_0}}^n\le M q^n ,
\]
и ряд мажорируется сходящейся геометрической прогрессией — абсолютная сходимость. Вторая часть —
контрапозиция первой.
\end{proof}

Следствие: область сходимости — это интервал $(-R,R)$ (возможно, с концами), где $R$ —
\emph{радиус сходимости}.

\section{Радиус сходимости}

\begin{definition}
\[
  R=\lim_{n\to\infty}\abs[\Big]{\frac{a_n}{a_{n+1}}}\quad(\text{по Даламберу, если предел есть}),
  \qquad
  R=\frac{1}{\varlimsup\limits_{n\to\infty}\sqrt[n]{\abs{a_n}}}\quad(\text{по Коши}).
\]
\end{definition}

\begin{theorem}[Коши–Адамар]
Ряд $\sum a_n x^n$ сходится абсолютно на $(-R,R)$ и расходится на
$(-\infty,-R)\cup(R,+\infty)$, где $R=1/\varlimsup\sqrt[n]{\abs{a_n}}$.
\end{theorem}

\begin{proof}
Применим радикальный признак Коши к $\sum\abs{a_n x^n}$:
$q(x)=\varlimsup\sqrt[n]{\abs{a_n x^n}}=\abs x\cdot\varlimsup\sqrt[n]{\abs{a_n}}=\dfrac{\abs x}{R}$.
При $q<1$, т.е. $\abs x<R$, ряд сходится; при $q>1$, т.е. $\abs x>R$, общий член не стремится к
нулю — расходимость. Случаи $R=0$ и $R=+\infty$ соответствуют $\varlimsup=+\infty$ и $0$.
\end{proof}

\begin{example}
$\displaystyle\sum_{n=1}^\infty\frac{x^n}{n}$: $\dfrac1R=\varlimsup\sqrt[n]{1/n}=1$, $R=1$. При
$x=1$ — гармонический ряд (расходится), при $x=-1$ — знакочередующийся (сходится). Область
сходимости $[-1,1)$.
\end{example}

\section{Свойства степенного ряда внутри интервала сходимости}

\begin{theorem}[равномерная сходимость]
$\sum a_n x^n$ сходится равномерно на любом отрезке $[-\rho,\rho]\subset(-R,R)$.
\end{theorem}

\begin{proof}
При $\abs x\le\rho$ имеем $\abs{a_n x^n}\le\abs{a_n}\rho^n$, а ряд $\sum\abs{a_n}\rho^n$ сходится
(так как $\rho<R$). По признаку Вейерштрасса сходимость равномерна.
\end{proof}

\begin{corollary}\leavevmode
\begin{enumerate}[label=\arabic*),leftmargin=*]
  \item \textbf{Непрерывность:} сумма $S(x)$ непрерывна на $(-R,R)$.
  \item \textbf{Почленное интегрирование:}
        $\displaystyle\int_0^x S(t)\dt=\sum_{n=0}^\infty\frac{a_n}{n+1}x^{n+1}$ с тем же радиусом
        сходимости (используется $\sqrt[n]{n}\to1$).
  \item \textbf{Почленное дифференцирование:}
        $\displaystyle S'(x)=\sum_{n=1}^\infty n\,a_n x^{n-1}$ с тем же радиусом сходимости.
\end{enumerate}
\end{corollary}

\begin{theorem}[вторая теорема Абеля]
Если $\sum a_n x^n$ сходится при $x=R$ (или $x=-R$), то $S(x)$ непрерывна в этой точке слева
(соответственно справа).
\end{theorem}

\begin{example}[вывод $\ln2$]
Из $\displaystyle\sum_{n=0}^\infty x^n=\frac1{1-x}$ заменой $x\mapsto-x$ и почленным
интегрированием:
\[
  \ln(1+x)=\sum_{n=1}^\infty\frac{(-1)^{n+1}}{n}x^n,\qquad \abs x<1 .
\]
При $x=1$ ряд $\sum\dfrac{(-1)^{n+1}}{n}$ сходится (Лейбниц), и по второй теореме Абеля
$\ln2=1-\tfrac12+\tfrac13-\tfrac14+\dots$.
\end{example}

\section{Суммирование степенных рядов}

Суммы находят, сводя ряд к геометрической прогрессии $\sum x^n=\dfrac1{1-x}$ и применяя почленное
дифференцирование/интегрирование.

\begin{example}
$\displaystyle\sum_{n=1}^\infty n x^n$ ($R=1$): дифференцируя $\sum_{n\ge0}x^n=\dfrac1{1-x}$,
получаем $\sum_{n\ge1}n x^{n-1}=\dfrac1{(1-x)^2}$, откуда
\[
  \sum_{n=1}^\infty n x^n=x\cdot\frac{1}{(1-x)^2}=\frac{x}{(1-x)^2}.
\]
\end{example}

\begin{example}
$\displaystyle\sum_{n=1}^\infty\frac{x^n}{n}$ ($R=1$): интегрируя $\sum_{n\ge0}x^n=\dfrac1{1-x}$,
получаем $\displaystyle\sum_{n=1}^\infty\frac{x^n}{n}=-\ln(1-x)$, $\abs x<1$.
\end{example}

\end{document}
